\scanpage{36}
\mspage{34}
% archive sheet number: 35
en remarquant que ce groupe algébrique sur $\unsure{\mathbb Q_\ell|Q_\ell}$ provient d'un groupe algébrique \deleted{sur} $\mathbb G_{a\,\mathbb Q}\times\Gamma_{\mathbb Q}$ sur $\mathbb Q$. On peut demander \inserted{d'abord} \deleted{que} si les représentations algébriques sur les différents $\mathbb Q_\ell$ définies par les $H^i(\overline X,\mathbb Q_\ell)$ proviennent d'une même représentation algébrique \uline{sur} $\mathbb Q$.
\deleted{(Ça devrait \unsure{sortir|rester} \unsure{un exercice|en exercice} \unsure{pour $i=1$}.)}
\inserted{Ce n'est raisonnable \deleted{que} cependant qu'en car. nulle, où une démonstration par voie transcendante (utilisant la cohomologie entière transcendante) devrait être possible. En car. résiduelle $>0$, on peut \unsure{vouloir|voir} que ces représentations s'écrivent (de façon unique à isom près) sur $\mathbb Q$; on peut demander qu'elles y soient indépendantes de $\ell$.}
% slanted
\margin{[N.B. \unsure{Même|M\textsuperscript{me}} la représentation induite de $\Gamma$ ne s'écrit pas en général sur $\mathbb Q$, si $i=1$ et si $X_K$ est une courbe elliptique sur $K$ ayant ``bonne réduction potentielle''.]}

Dans le cas d'un corps résiduel fini, on trouve un groupe ``loc. algébrique''
\[
 \mathbb Z_{\mathbb Q}\cdot(\mathbb G_{a\,\mathbb Q}\times\Gamma_{\mathbb Q})=G
\]
sur $\mathbb Q$, produit semi-direct, où \inserted{le générateur} $f_q\in\mathbb Z_{\mathbb Q}(\mathbb Q)$ opère sur $\mathbb G_{a\,\mathbb Q}$ \deleted{\unsure{par une}} par multiplication par $q$, et sur $\Gamma$ comme il faut [sur $\unsure{(\Gamma/[\Gamma,\Gamma])(\ell)|(\Gamma/[\Gamma,\Gamma])^{(\ell)}|(\Gamma/[\Gamma,\Gamma])_\ell}$ il opère aussi par multiplication par $q$]. On peut encore demander que les représentations sur les différents $\mathbb Q_\ell$ ($\ell\ne p$) définies par les $H^i(\overline X,\mathbb Q_\ell)$ soient \deleted{définies sur $\mathbb Q$, et déduites} \inserted{isomorphes à celles déduites} d'une même représentation de $G$, sur $\mathbb Q$.
